% Planned 210.8; source-native IEL continuation, not a PMCE extension.
\section{Question-to-question implication}
\label{sec:p18-question-to-question-implication}

An inquiry can warrant a further question without already having its answer.
The relevant guarantee is not that the first question is true and therefore
the second is true: questions are not declarative formulas here. It is that
an initially sound question, together with true auxiliary information, leads
to a sound further question, and that every possible direct answer to the
further question helps restrict the initial alternatives. These are separate
requirements. An unrelated yes--no question meets the first requirement in
classical logic but need not meet the second.

Wi\'sniewski's inferential erotetic logic (IEL) supplies the two-clause relation
\cite[secs.~3--4,7]{Wisniewski2001QuestionsInferences}. His later deduction and
reduction results are formulated in minimal erotetic semantics rather than
only the finite propositional calculus \cite{Wisniewski2018DeductionReduction}.
The present treatment develops that relation, its source conditions, complete
semantic proof routes, and a worked inquiry. The preceding PMCE chapter
remains a calculus for declarative-premise evocation, not a proof system for
the relation defined below.

\subsection{Language, admissible situations and answerhood}

\begin{hbdefinition}{def:p18-ei-setting}{The minimal semantic setting}
Fix a language \(L\) with disjoint categories of declarative formulas
\(D_L\) and questions. A syntactic function \(d\) assigns to each question
\(Q\) a set \(dQ\subseteq D_L\) with at least two elements. The set may
be infinite. In the first-order instances below its members are closed
sentences, not formulas with an unassigned individual variable. Being a
direct answer does not assert its truth or its uniqueness.

Fix a class \(\mathcal P\) of admissible partitions
\(P=(T_P,U_P)\) of \(D_L\): the two sets are disjoint and have union
\(D_L\). Membership in \(T_P\) is the designated truth notion. The class
must be specified by the chosen declarative semantics; it is not the class
of all arbitrary assignments to complex formulas. A question is
\emph{sound at} \(P\) exactly when \(dQ\cap T_P\ne\varnothing\).
Define
\[
 X\models_L A\quad\Longleftrightarrow\quad
 (\forall P\in\mathcal P)(X\subseteq T_P\Rightarrow A\in T_P),
\]
\[
 X\mathrel{\|=}_L Y\quad\Longleftrightarrow\quad
 (\forall P\in\mathcal P)(X\subseteq T_P\Rightarrow Y\cap T_P\ne\varnothing).
\]
The second relation is multiple-conclusion entailment, not entailment of
every member of \(Y\). The double-bar symbol has the same local role as
in the preceding chapter, now with the underlying language displayed.
Write \(A\models_L C\) for \(\{A\}\models_L C\), and use commas
in premise positions for union with the displayed declaratives.
\end{hbdefinition}

Boolean valuations give one admissible class. First-order models on nonempty
domains give another, with the interpretation of constants and equality fixed
in each model. Restricting that class to models in which every individual has
a name is a further assumption, not a consequence of writing a wh-question.
IEL does not require its declarative basis to be classical in every application.

\begin{hbdefinition}{def:p18-ei-consequence-setting}{The source's consequence-generated partitions}
The 2018 paper provides a particular way to specify \(\mathcal P\). Fix a
nonempty, nonuniversal consequence relation \(\vdash_\ell\) with overlap,
monotonicity in the premises, and cut for sets. Write
\(\operatorname{Cn}_\ell(Y)=\{A:Y\vdash_\ell A\}\).
A set \(Y\) is maximally consistent in the source's sense when
\(\operatorname{Cn}_\ell(Y)\ne D_L\) and, for every \(A\), either
\(A\in\operatorname{Cn}_\ell(Y)\) or
\(\operatorname{Cn}_\ell(Y\cup\{A\})=D_L\).
Assume the \emph{Lb feature}: whenever \(X\not\vdash_\ell A\), there
exists such a \(Y\supseteq X\) with \(Y\not\vdash_\ell A\).
Then take all partitions whose true component is
\(\operatorname{Cn}_\ell(Y)\) for such a \(Y\).
Lb is an assumption on the consequence relation, not a theorem for every
nonclassical logic \cite[pp.~298--300]{Wisniewski2018DeductionReduction}.
\end{hbdefinition}

\begin{hbtheorem}{thm:p18-ei-basis}{proved=annex:p18.ei.basis}{Adequacy of the consequence-generated semantic basis}
Under Definition~\ref{def:p18-ei-consequence-setting},
\(X\models_L A\) if and only if \(X\vdash_\ell A\).
This adequacy is for the declarative basis; it does not assert completeness
of a syntactic calculus for question-to-question implication.
\end{hbtheorem}

\subsection{The two clauses, and what they do not say}

\begin{hbdefinition}{def:p18-ei-relation}{Erotetic implication and its regular special case}
For comparison, the general semantic evocation relation is
\[
 E_L(X,Q)\quad\Longleftrightarrow\quad
 X\mathrel{\|=}_L dQ\ \text{and}\
 (\forall A\in dQ)\ X\not\models_L A.
\]
It has only declarative premises. The following relation instead has a
question premise as well as its declarative basis.
The question \(Q\) \emph{erotetically implies} \(R\) on basis \(X\),
written \(\operatorname{Im}_L(Q,X,R)\), when both conditions hold:
\begin{enumerate}
\item For every \(A\in dQ\), \(X\cup\{A\}\mathrel{\|=}_L dR\).
\item For every \(B\in dR\), there is a set \(Y_B\) with
\(\varnothing\ne Y_B\subsetneq dQ\) and
\(X\cup\{B\}\mathrel{\|=}_L Y_B\).
\end{enumerate}
The witnesses are answer sets, not necessarily singleton answers. Pure
implication means \(X=\varnothing\). \emph{Regular} implication,
written here \(\operatorname{RIm}_L(Q,X,R)\), retains the first condition
and replaces the second by: for every \(B\in dR\), some \(A\in dQ\)
satisfies \(X\cup\{B\}\models_L A\). The abbreviation RIm is editorial;
the source calls this regular erotetic implication
\cite[pp.~23--26]{Wisniewski2001QuestionsInferences}.
\end{hbdefinition}

\begin{hbtheorem}{thm:p18-ei-profile}{proved=annex:p18.ei.profile}{Soundness transmission, monotonicity and boundaries}
Erotetic implication transmits soundness in both directions at any admissible
partition making all members of \(X\) true. Regular implication implies
erotetic implication. Both relations are preserved on enlarging \(X\).
If no admissible partition satisfies \(X\), both relations hold vacuously
for any well-formed questions. This contrasts with evocation
\(E_L(X,Q)\), which additionally requires that no direct answer be entailed.

General implication together with \(E_L(X,Q)\) need not yield
\(E_L(X,R)\). Regular implication together with \(E_L(X,Q)\) does yield it.
Multiple true direct answers are allowed; a witness \(Y_B\) guarantees a
true answer in that subset, not the falsity of all answers outside it.
\end{hbtheorem}

For example, take \(Q=?\{p,q,r\}\), \(R=?\{p\lor q,r\}\), and
\(X=\{p\lor q\}\) in the preceding Boolean language. The witnesses
for \(R\)'s answers are \(\{p,q\}\) and \(\{r\}\), so implication
holds. The initial question is evoked, but \(p\lor q\) already settles
one answer to \(R\), so \(R\) is not evoked. At \(p=q=r=1\), accepting
\(p\lor q\) does not eliminate \(r\). A reduction in the class containing
a guaranteed answer is not necessarily elimination of every other true answer.

An empty direct-answer set or a singleton is excluded by the source language.
A \emph{null factual situation}, where none of a well-formed question's
answers is true, is different: the question is unsound there. With true basis
\(X\), implication makes the two questions sound or unsound together. A
negative answer such as ``none of these'' can be added as a declarative, but
doing so changes the question and requires a new implication check. In
particular, empty \(X\) means pure implication, not an empty question.
For multiple-conclusion entailment, \(X\mathrel{\|=}_L\varnothing\) holds
exactly when \(X\) has no admissible satisfying partition; allowing that
empty set as a reduction witness would erase the intended answer restriction.

\subsection{Presupposition is not implication}

\begin{hbdefinition}{def:p18-ei-presupposition}{Presuppositions and partial answers}
\(C\) is a presupposition of \(Q\) when \(A\models_L C\) for every
\(A\in dQ\); write the set of these as \(\operatorname{Pres}(Q)\).
It is prospective when also \(\{C\}\mathrel{\|=}_L dQ\).
A question is normal when
\(\operatorname{Pres}(Q)\mathrel{\|=}_L dQ\), and regular when it has a
prospective presupposition. Regularity of a \emph{question} is not regularity
of an \emph{implication}. A just-complete answer is equivalent under
\(\models_L\) to a direct answer. A partial answer \(C\) is not
just-complete but has a nonempty proper witness \(Y\subsetneq dQ\) with
\(\{C\}\mathrel{\|=}_L Y\) and \(A\models_L C\) for all \(A\in Y\)
\cite[pp.~15--17]{Wisniewski2001QuestionsInferences}.
\end{hbdefinition}

\begin{hbtheorem}{thm:p18-ei-presuppositions}{proved=annex:p18.ei.presuppositions}{What presupposition transmission establishes}
Evocation entails every presupposition of the evoked question. For normal
\(Q\), evocation is equivalent to entailing all its presuppositions and no
direct answer. For regular \(Q\), it is equivalent to entailing some
prospective presupposition and no direct answer.

If \(Q,R\) are normal, the first clause of erotetic implication alone is
equivalent to
\[
 (\forall C\in\operatorname{Pres}(R))\quad
 X\cup\operatorname{Pres}(Q)\models_L C.
\]
If both are regular questions, that first clause is equivalent to the existence
of prospective presuppositions \(D\) of \(Q\) and \(C\) of \(R\) with
\(X\cup\{D\}\models_L C\). Neither characterization supplies the
second, answer-reduction clause. General semantics need not make normality
equivalent to regularity; compactness is not assumed here.
Under the stated partial-answer definition, a question with exactly two
direct answers has no partial answers. This does not say that it lacks
incomplete or corrective responses of other kinds.
\end{hbtheorem}

In the finite Boolean language, \(p\lor q\lor r\) is a prospective
presupposition of \(?\{p,q,r\}\). Asking the independent polar question
\(?\{s,\neg s\}\) preserves soundness and its tautological presupposition.
But neither answer points to any proper subset of \(\{p,q,r\}\): set
only an excluded original alternative true, with the selected value of \(s\).
Thus presupposition transmission alone is not a valid inquiry reduction.

\subsection{Internalizing declarative premises}

\begin{hbdefinition}{def:p18-ei-conditional-language}{The conditional language conditions}
For these results assume a declarative connective \(\to\) satisfying, at
every admissible partition,
\[
 (C\to A)\in T_P\quad\Longleftrightarrow\quad
 C\notin T_P\ \text{or}\ A\in T_P.
\]
Assume the source's question closure conditions: for every question \(Q\)
and declarative \(C\), there exists \(Q^C\) with
\(dQ^C=\{C\to A:A\in dQ\}\); and a question with such a prefixed
answer set has an unprefixed counterpart. The syntactic prefix map is
injective on formula expressions. Define \(F_C(Y)=\{C\to A:A\in Y\}\).
This denotes a set of declaratives, not a conditional with a question as its
argument. The truth clause is a real restriction: the source explicitly does
not infer it for every nonclassical implication
\cite[pp.~301--302]{Wisniewski2018DeductionReduction}.
\end{hbdefinition}

\begin{hblemma}{thm:p18-ei-conditional-lemmas}{proved=annex:p18.ei.conditional-lemmas}{The three conditional entailment transformations}
For \(Y\ne\varnothing\), under the conditional language conditions,
\[
 X,C\mathrel{\|=}_L Y\ \Longleftrightarrow\ X\mathrel{\|=}_L F_C(Y),
\]
\[
 X,D\mathrel{\|=}_L F_C(Y)\ \Longleftrightarrow\
 X,C\to D\mathrel{\|=}_L F_C(Y),
\]
\[
 X,C,D\mathrel{\|=}_L Y\ \Longleftrightarrow\
 X,C\to D\mathrel{\|=}_L F_C(Y).
\]
Comma means set union with the displayed declaratives. The nonempty-set
condition cannot simply be deleted from these source lemmas.
\end{hblemma}

\begin{hbtheorem}{thm:p18-ei-evocation-deduction}{proved=annex:p18.ei.evocation-deduction}{The general evocation bridge}
\(E_L(X\cup\{C\},Q)\) if and only if \(E_L(X,Q^C)\).
For finite \(X=\{C_1,\ldots,C_n\}\), prefix every answer by the nested
conditional \(C_1\to(C_2\to\cdots\to(C_n\to A)\cdots)\).
Then evocation from \(X\) is equivalent to evocation of that transformed
question from the empty set. For \(n=0\) use the identity transformation.
This is the semantic generalization of the preceding finite PMCE bridge,
not first-order completeness of PMCE.
\end{hbtheorem}

\begin{hbtheorem}{thm:p18-ei-deduction}{proved=annex:p18.ei.deduction}{Deduction and finite reduction for question implication}
Under Definition~\ref{def:p18-ei-conditional-language},
\[
 \operatorname{Im}_L(Q,X\cup\{C\},R)
 \quad\Longleftrightarrow\quad
 \operatorname{Im}_L(Q^C,X,R^C).
\]
Both questions, not only the conclusion, must be transformed. If \(X\) is
finite, repeated conditionalization gives
\[
 \operatorname{Im}_L(Q,X,R)
 \quad\Longleftrightarrow\quad
 \operatorname{Im}_L(Q^{\vec C},\varnothing,R^{\vec C}),
\]
where \(\vec C\) lists all elements of \(X\) without repetitions, and
both direct-answer sets receive the same nested prefix in that order.
No infinite sequence of prefixes is asserted for an infinite basis.
\end{hbtheorem}

If \(C\) is false, all the \(C\)-conditionalized answers are true. This
does not spoil the implication theorem: the two questions are still
syntactically well formed, and the same injective map preserves their proper
subsets. Semantic equivalence of all the answers at that particular partition
is not syntactic identity. Evocation, by contrast, still has its separate
non-entailment requirements.

\begin{hbdefinition}{def:p18-ei-conjunct-language}{The conjunctive alternative and exact syntax scope}
Assume instead the truth clause
\((C\land A)\in T_P\) iff \(C,A\in T_P\).
For a finite nonempty \(X=\{C_1,\ldots,C_n\}\), let
\(G_X(A)=C_1\land(C_2\land\cdots\land(C_n\land A)\cdots)\).
Extend this map to sets by \(G_X(Y)=\{G_X(A):A\in Y\}\).
For the following nested version explicitly require that questions
\(Q_X,R_X\) with answer sets \(G_X(dQ),G_X(dR)\) exist; injectivity
of the syntactic prefix is also required. The finite explicit question
language supplies this condition. A different language must establish it.

The source calls a question added-conjunct when its answers have a uniform
conjunct prefix derived from a question without a common conjunct form.
Its printed closure CQ1-conjunction is restricted to a question which is
not already added-conjunct. That restricted clause alone does not license
successive applications to nested added-conjunct questions. Alternatively,
form a single bundled premise \(C_X=C_1\land\cdots\land C_n\) and
prefix answers once by \(C_X\land A\): for original non-added-conjunct
questions the source's restricted closure licenses this single operation.
The nested and bundled expressions are truth equivalent under the stated
clause, but existence of one question expression does not define the other
\cite[pp.~305--307]{Wisniewski2018DeductionReduction}.
\end{hbdefinition}

\begin{hblemma}{thm:p18-ei-conjunct-syntax-boundary}{proved=annex:p18.ei.conjunct-syntax-boundary}{Why restricted one-step closure does not provide nested questions}
There is a Boolean question language satisfying the source's restricted
CQ1-conjunction and inverse CQ2-conjunction conditions, but lacking a
question with direct answers \(\{a\land(b\land p),a\land(b\land q)\}\),
although it has a non-added-conjunct question with answers \(\{p,q\}\).
The variables are syntactically distinct. Thus the extra exact-question
existence condition in the nested reduction is substantive; it is not
established merely by associativity of conjunction.
\end{hblemma}

\begin{hbtheorem}{thm:p18-ei-conjunct-reduction}{proved=annex:p18.ei.conjunct-reduction}{Conjunctive reduction with explicit question existence}
Under Definition~\ref{def:p18-ei-conjunct-language},
\(\operatorname{Im}_L(Q,X,R)\) if and only if
\(\operatorname{Im}_L(Q_X,\varnothing,R_X)\).
The same equivalence holds for the one-bundled-conjunct variant when those
questions exist. This alternative needs the conjunction truth clause, not
the conditional truth clause. Empty \(X\) uses the identity, not a new
primitive empty conjunction.
\end{hbtheorem}

\subsection{Composition: an actual restriction, not unrestricted transitivity}

\begin{hbtheorem}{thm:p18-ei-composition}{proved=annex:p18.ei.composition}{Answer-reduction composition}
If \(\operatorname{Im}_L(Q,X,R)\) and
\(\operatorname{RIm}_L(R,Z,S)\), then
\(\operatorname{Im}_L(Q,X\cup Z,S)\).
If both legs are regular, the conclusion is regular.

More generally, supply explicit witnesses \(Y_B\subsetneq dQ\) for the
first leg and \(V_C\subsetneq dR\) for the second. Composition is valid
if, for every \(C\in dS\),
\(\bigcup_{B\in V_C}Y_B\) is a proper subset of \(dQ\).
All witnesses must be nonempty. Without this last properness condition,
ordinary erotetic implication need not compose.
\end{hbtheorem}

\begin{hbtheorem}{thm:p18-ei-nontransitivity}{proved=annex:p18.ei.nontransitivity}{A complete countercase to unrestricted composition}
In the Boolean language let
\[
 Q=?\{p,q,r\},\quad
 R=?\{p\lor q,p\lor r,q\lor r\},\quad
 S=?\{p\lor q\lor r,s\},\quad X=\{s\to p\}.
\]
Then \(\operatorname{Im}(Q,X,R)\) and
\(\operatorname{Im}(R,X,S)\), but not
\(\operatorname{Im}(Q,X,S)\). All displayed direct answers are satisfiable;
the failure does not rely on inserting a contradictory target answer.
\end{hbtheorem}

The rule sheet records semantic certification and the proved transformations.
It is not an asserted complete analytic calculus: its certification entry
requires actual declarative multiple-conclusion entailments and explicit
answer-set witnesses. A source theorem's name is not itself a certificate.
\calculussheet{sheet:p18-erotetic-implication}{IEL semantic implication certificates and proved transformations}
  {A fixed language with declaratives, questions, at least two direct answers
  per question, and a fixed admissible partition class. These are semantic
  certification/transform entries, not primitive rules of PMCE and not a
  claimed complete analytic calculus. Definitions~\ref{def:p18-ei-setting}
  and~\ref{def:p18-ei-relation} fix the entailment and witness notation.
  Conditional and conjunctive entries require the separate displayed language
  hypotheses. All questions in an entry must actually exist and be well formed.}

\hbrule{rule:p18-ei-certify}{EI certificate (definitional)}
  {(\forall A\in dQ)\ X,A\mathrel{\|=}_L dR\qquad
   (\forall B\in dR)\ X,B\mathrel{\|=}_L Y_B}
  {\operatorname{Im}_L(Q,X,R)}
  {For every \(B\), an explicit witness
   \(\varnothing\ne Y_B\subsetneq dQ\) is supplied. Each displayed
   declarative entailment must be established, not supplied by a theorem name.}

\hbrule{rule:p18-ei-regular-certify}{Regular certificate (definitional)}
  {(\forall A\in dQ)\ X,A\mathrel{\|=}_L dR\qquad
   (\forall B\in dR)\ X,B\models_L a_B}
  {\operatorname{RIm}_L(Q,X,R)}
  {For every \(B\), \(a_B\in dQ\) is specified. The singleton
   \(\{a_B\}\) is proper because \(|dQ|\geq2\). RIm is the
   editorial abbreviation for the source's regular implication.}

\hbrule{rule:p18-ei-basis-extension}{Basis extension (proved)}
  {\operatorname{Im}_L(Q,X,R)}
  {\operatorname{Im}_L(Q,X\cup Z,R)}
  {\(Z\subseteq D_L\); same language and partition class.
   Existing witness sets are retained. The same entry holds for RIm;
   it is not a weakening rule for evocation. Proof: \ref{annex:p18.ei.profile}.}

\hbrule{rule:p18-ei-regular-compose}{Regular-last-leg composition (proved)}
  {\operatorname{Im}_L(Q,X,R)\qquad\operatorname{RIm}_L(R,Z,S)}
  {\operatorname{Im}_L(Q,X\cup Z,S)}
  {Same language, same admissible class and same intermediate direct-answer
   set. If the first leg is regular, the conclusion is regular.
   Proof: \ref{annex:p18.ei.composition}.}

\hbrule{rule:p18-ei-witness-compose}{Explicit-witness composition (proved)}
  {\operatorname{Im}_L(Q,X,R)\qquad\operatorname{Im}_L(R,Z,S)}
  {\operatorname{Im}_L(Q,X\cup Z,S)}
  {The first leg has specified witnesses \(Y_B\) and the second
   specified witnesses \(V_C\). For every \(C\in dS\), require
   \(\bigcup_{B\in V_C}Y_B\subsetneq dQ\).
   All individual witness sets are nonempty. Without this extra properness
   condition the entry is invalid. Proof: \ref{annex:p18.ei.composition}.}

\hbrule{rule:p18-ei-condition-remove}{Conditional premise removal (proved)}
  {\operatorname{Im}_L(Q,X\cup\{C\},R)}
  {\operatorname{Im}_L(Q^C,X,R^C)}
  {The conditional truth clause and question closure of
   Definition~\ref{def:p18-ei-conditional-language};
   \(dQ^C=F_C(dQ)\), \(dR^C=F_C(dR)\), with injective syntactic
   prefix. Both questions are transformed. Proof: \ref{annex:p18.ei.deduction}.}

\hbrule{rule:p18-ei-condition-restore}{Conditional premise restoration (proved)}
  {\operatorname{Im}_L(Q^C,X,R^C)}
  {\operatorname{Im}_L(Q,X\cup\{C\},R)}
  {The same conditional truth clause; the unprefixed questions exist by
   inverse closure or are already given. Their answer sets have exactly the
   stated prefix images, whose inverse images preserve nonempty proper witnesses.
   Proof: \ref{annex:p18.ei.deduction}.}

\hbrule{rule:p18-ei-conjunct-remove}{Finite conjunctive premise removal (proved)}
  {\operatorname{Im}_L(Q,X,R)}
  {\operatorname{Im}_L(Q_X,\varnothing,R_X)}
  {\(X\) finite nonempty; conjunction truth clause, injective prefix and
   actual existence of both exact transformed questions as in
   Definition~\ref{def:p18-ei-conjunct-language}. Restricted closure for
   non-added-conjunct questions is not silently iterated.
   Proof: \ref{annex:p18.ei.conjunct-reduction}.}

\hbrule{rule:p18-ei-conjunct-restore}{Finite conjunctive premise restoration (proved)}
  {\operatorname{Im}_L(Q_X,\varnothing,R_X)}
  {\operatorname{Im}_L(Q,X,R)}
  {The same finite prefix, truth clause and exactly corresponding original
   and transformed questions; use inverse images of witnesses. The bundled
   one-conjunct variant has the separately stated existence condition.
   Proof: \ref{annex:p18.ei.conjunct-reduction}.}


\subsection{Two source-native implication patterns}

\begin{hblemma}{thm:p18-ei-native-patterns}{proved=annex:p18.ei.native-patterns}{Refinement and auxiliary equivalence}
For distinct Boolean variables \(p,q,r\), let
\(Q=?\{p,q\lor r\}\) and \(R=?\{p,q,r\}\).
Then \(Q\) purely and regularly implies \(R\); \(R\) purely implies
\(Q\), but that reverse implication is not regular. These instantiate
the source's refinement/coarsening examples (8.4)--(8.5).

For well-formed finite questions \(?\{A_1,\ldots,A_n\}\) and
\(?\{B_1,\ldots,B_n\}\), \(n>1\), the basis
\(X=\{A_i\leftrightarrow B_i:1\leq i\leq n\}\) gives regular
implication in both directions under Boolean semantics. This is source
pattern (8.17), with every corresponding answer and every auxiliary
biconditional kept explicit \cite[pp.~27--28]{Wisniewski2001QuestionsInferences}.
\end{hblemma}

Refinement and coarsening have the same coverage condition, but their
answer-reduction strengths differ. The coarser answer \(q\lor r\) need
not settle either finer answer. Auxiliary equivalence is different again:
each corresponding answer is transferred individually, but only where its
auxiliary basis is true. For example, put
\(Q=?\{p,r\}\), \(R=?\{q,s\}\), and
\(X=\{p\leftrightarrow q,r\leftrightarrow s\}\).
The relation holds. At \((p,r,q,s)=(1,0,0,0)\), however, the auxiliary
basis is false, \(Q\) is sound, and \(R\) is unsound. At
\((0,0,1,0)\) the reverse soundness mismatch occurs. These valuations
do not refute the relation; they refute a claim that its guarantee is
unconditional on the declarative warrants.

\subsection{A worked inquiry with a genuine partial-answer stage}

\begin{hbexample}{def:p18-ei-travel-inquiry}{Destination, departure time and a branch question}
Use the source's Andrew example, construed here as a formal inquiry rather
than an empirical travel report \cite[p.~25]{Wisniewski2001QuestionsInferences}.
Let \(p,l,m\) mean departures for Paris, London, Moscow; let \(a,e\)
mean morning and evening departure. The initial question is
\(Q=?\{p,l,m\}\), the first subquestion \(R=?\{a,e\}\), and the
auxiliary basis is
\[
 X=\{(p\lor l\lor m)\to(a\lor e),\ a\to(p\lor l),\ e\to m\}.
\]
Every initial direct answer guarantees an answer to \(R\).
The witnesses are \(Y_a=\{p,l\}\) and \(Y_e=\{m\}\).
A morning answer restricts the original alternatives but does not yet
establish Paris or London separately; an evening answer establishes Moscow.
No premise asserts exclusive destinations, so the proof does not infer
\(\neg m\) from a morning answer.

Suppose the accepted answer is \(a\). Update the declarative basis to
\(Z=X\cup\{a\}\). It covers the branch question \(B=?\{p,l\}\)
but settles neither answer. The polar subquestion \(T=?\{p,\neg p\}\)
regularly follows from \(B\) on \(Z\): a positive response yields \(p\);
a negative response, together with \(p\lor l\), yields \(l\).
The inquiry therefore finds a direct answer to the original question on
either response. This is a conditional branch with an updated basis, not
an unproved globally transitive chain on the old \(X\).
\end{hbexample}

\begin{hbtheorem}{thm:p18-ei-inquiry}{proved=annex:p18.ei.inquiry}{Every branch and a failed warrant}
The stated travel inquiry has all the implication, evocation and resolution
properties in Definition~\ref{def:p18-ei-travel-inquiry}. Its evening branch
resolves \(m\); its morning branch followed by either truthful polar response
resolves \(p\) or \(l\). If \(a\to(p\lor l)\) is omitted, the first
question-to-question implication fails: a morning answer can coexist with
all three original answers false. The accepted answer's truth and the
auxiliary warrants are essential to actual inquiry success.
\end{hbtheorem}

\subsection{Named witnesses, complete lists and domain failures}

\begin{hbdefinition}{def:p18-ei-wh-language}{The source's monadic wh-question instance}
Use monadic first-order logic with equality, individual constants, no
function symbols, and nonempty model domains. Here \(Ax\) has \(x\) as
its only free variable; substituting a constant gives a closed sentence.
Besides finite explicit
questions, the 2001 source introduces \(?S(Ax)\), whose direct answers
are all closed instances \(A(x/c)\) for constants \(c\), and
\(?U(Ax)\), whose direct answers, for finite nonempty lists of distinct
constant symbols \(c_1,\ldots,c_k\), are
\[
 A(x/c_1)\land\cdots\land A(x/c_k)\land
 \forall x\bigl(Ax\to(x=c_1\lor\cdots\lor x=c_k)\bigr).
\]
The former asks for a named witness, the latter for a complete finite list.
There are enough constant symbols that each direct-answer set has at least
two syntactically distinct sentences. Distinct symbols need not denote
distinct individuals: no unique-names assumption is introduced
\cite[pp.~20--22,29]{Wisniewski2001QuestionsInferences}.
\end{hbdefinition}

\begin{hbtheorem}{thm:p18-ei-wh-cases}{proved=annex:p18.ei.wh-cases}{Regular wh-reduction and its soundness limits}
On \(X=\{\forall x(Ax\leftrightarrow Bx)\}\), \(?S(Ax)\) regularly
implies \(?S(Bx)\), and \(?U(Ax)\) regularly implies \(?U(Bx)\),
over the stated class of all models. On
\(\{\forall x(Ax\to(x=c_1\lor\cdots\lor x=c_n))\}\), with \(n>1\),
\(?S(Ax)\) regularly implies \(?\{A(x/c_1),\ldots,A(x/c_n)\}\).

Neither \(\exists xAx\) nor nonempty domain alone guarantees soundness
of \(?S(Ax)\) over arbitrary models: the witness may be unnamed.
\(?U(Ax)\) can be unsound when the extension is empty, when it is infinite,
or when a required individual lacks a name. Even a fully named nonempty
domain does not make an empty extension a nonempty complete-list answer.
The source itself cautions that more general complete-list analyses need
additional quantificational resources; no general wh-completeness calculus
is asserted here.
\end{hbtheorem}

\subsection{Finite certificates and standing of the continuation}

\begin{hbmethod}{def:p18-ei-certificate}{A finite Boolean certificate}
For finite explicit answer sets and a finite Boolean basis, supply for every
target answer \(B\) a nonempty proper index subset of the source answers.
Check every assignment to the finitely many occurring variables. At any
assignment satisfying the basis and some source answer, check that a target
answer is true. At any assignment satisfying the basis and \(B\), check
that an answer in its chosen subset is true. For regular certification choose
singleton subsets. Question well-formedness is checked separately.
\end{hbmethod}

\begin{hbtheorem}{thm:p18-ei-certificate}{proved=annex:p18.ei.certificate}{Exactness and termination of finite certification}
The certificate is sound and complete for the specified finite Boolean
implication relation: some such witness certificate passes if and only if
the relation holds. Exhaustive search over the finitely many nonempty proper
subsets terminates, as does assignment checking. This is an editorial
semantic procedure, not a claim that PMCE proves question-premise sequents,
not an efficiency bound for general first-order inference, and not a substitute
for the general source semantics.
\end{hbtheorem}

The continuation supplies the actual IEL relation, the general declarative
semantic basis, source deduction/reduction constructions and their assumptions,
restricted composition, and worked propositional and wh-cases. It does not
identify IEL with Harrah's question logic, Ginzburg's context-update theory,
or Hintikka--Halonen--Mutanen's interrogative tableau framework. Those exact
systems are not identified by a translation or equivalence theorem here.
The claims above concern the specified language, answers and admissible
partition class. The 2018 article's historical statement about then-existing
axiomatizations is not asserted as the state of the literature today.
